\chapter{Foundations: Mathematics, Statistics, and Programming}
\label{ch:foundations}

This chapter establishes the mathematical, statistical, and programming
foundations needed to understand the rest of the book. If you are already
comfortable with these topics, feel free to skip ahead and use this chapter
as a reference.

\section{Why These Foundations Matter}

Quantitative finance sits at the intersection of three disciplines:

\begin{itemize}
    \item \textbf{Mathematics}: The language of precision---formulas, proofs,
      and rigorous reasoning
    \item \textbf{Statistics}: The science of uncertainty---data analysis,
      probability, and inference
    \item \textbf{Programming}: The tool of implementation---turning ideas
      into working systems
\end{itemize}

\marginnote{%
\textbf{The Quant Mindset}\\[0.3em]
A quantitative approach means making decisions based on data and models,
not intuition or emotion.
}

You don't need to be an expert in all three to start. This book teaches by
doing---each concept is introduced when needed and immediately applied.

\section{Mathematical Notation}

Throughout this book, we use standard mathematical notation. Here's a quick
reference:

\subsection{Sets and Numbers}

\begin{center}
\begin{tabular}{cl}
\toprule
\textbf{Symbol} & \textbf{Meaning} \\
\midrule
$\mathbb{R}$ & Real numbers (e.g., $3.14$, $-2.5$, $\sqrt{2}$) \\
$\mathbb{R}^+$ & Positive real numbers \\
$\mathbb{Z}$ & Integers ($\ldots, -2, -1, 0, 1, 2, \ldots$) \\
$\mathbb{N}$ & Natural numbers ($1, 2, 3, \ldots$) \\
$\in$ & ``is an element of'' (e.g., $x \in \mathbb{R}$ means $x$ is real) \\
$\forall$ & ``for all'' \\
$\exists$ & ``there exists'' \\
\bottomrule
\end{tabular}
\end{center}

\subsection{Summation and Products}

\marginnote{%
\textbf{Summation}\\[0.5em]
$\displaystyle\sum_{i=1}^{n} x_i = x_1 + x_2 + \cdots + x_n$
}

The summation symbol $\Sigma$ means ``add up'':
\[
\sum_{i=1}^{n} x_i = x_1 + x_2 + \cdots + x_n
\]

For example, if $x = [2, 5, 3]$, then $\sum_{i=1}^{3} x_i = 2 + 5 + 3 = 10$.

\marginnote{%
\textbf{Product}\\[0.5em]
$\displaystyle\prod_{i=1}^{n} x_i = x_1 \times x_2 \times \cdots \times x_n$
}

The product symbol $\Pi$ means ``multiply together'':
\[
\prod_{i=1}^{n} x_i = x_1 \times x_2 \times \cdots \times x_n
\]

\subsection{Functions and Calculus}

\begin{center}
\begin{tabular}{cl}
\toprule
\textbf{Symbol} & \textbf{Meaning} \\
\midrule
$f(x)$ & Function $f$ evaluated at $x$ \\
$f'(x)$ or $\frac{df}{dx}$ & Derivative of $f$ with respect to $x$ \\
$\frac{\partial f}{\partial x}$ & Partial derivative (when $f$ has multiple variables) \\
$\ln(x)$ & Natural logarithm (base $e \approx 2.718$) \\
$\exp(x)$ or $e^x$ & Exponential function \\
$\argmax_x f(x)$ & Value of $x$ that maximizes $f$ \\
$\argmin_x f(x)$ & Value of $x$ that minimizes $f$ \\
\bottomrule
\end{tabular}
\end{center}

\begin{keyinsight}
The natural logarithm $\ln$ is everywhere in finance. Key property:
$\ln(a \times b) = \ln(a) + \ln(b)$. This turns multiplication (compounding)
into addition, making calculations much easier.
\end{keyinsight}

\section{Statistics Essentials}

Statistics helps us understand and quantify uncertainty---essential when
dealing with financial markets.

\subsection{Measures of Central Tendency}

\marginnote{%
\textbf{Mean}\\[0.5em]
$\bar{x} = \dfrac{1}{n}\displaystyle\sum_{i=1}^{n} x_i$
}

The \textbf{mean} (average) is the sum divided by the count:
\[
\bar{x} = \frac{1}{n}\sum_{i=1}^{n} x_i
\]

The \textbf{median} is the middle value when data is sorted. It's more robust
to outliers than the mean.

\subsection{Measures of Spread}

\marginnote{%
\textbf{Variance}\\[0.5em]
$\sigma^2 = \dfrac{1}{n}\displaystyle\sum_{i=1}^{n} (x_i - \bar{x})^2$
}

\textbf{Variance} measures how spread out the data is:
\[
\sigma^2 = \frac{1}{n}\sum_{i=1}^{n} (x_i - \bar{x})^2
\]

\textbf{Standard deviation} is the square root of variance:
\[
\sigma = \sqrt{\sigma^2}
\]

Standard deviation has the same units as the original data, making it easier
to interpret than variance.

\begin{lstlisting}[language=Rust]
/// Calculate mean of a slice
pub fn mean(data: &[f64]) -> f64 {
    if data.is_empty() { return 0.0; }
    data.iter().sum::<f64>() / data.len() as f64
}

/// Calculate variance of a slice
pub fn variance(data: &[f64]) -> f64 {
    if data.is_empty() { return 0.0; }
    let m = mean(data);
    data.iter().map(|x| (x - m).powi(2)).sum::<f64>() / data.len() as f64
}

/// Calculate standard deviation
pub fn std_dev(data: &[f64]) -> f64 {
    variance(data).sqrt()
}
\end{lstlisting}

\subsection{Probability Distributions}

A \textbf{probability distribution} describes how likely different outcomes are.

\subsubsection{Normal Distribution (Gaussian)}

\marginnote{%
\textbf{Normal PDF}\\[0.5em]
$f(x) = \dfrac{1}{\sigma\sqrt{2\pi}} e^{-\frac{(x-\mu)^2}{2\sigma^2}}$
}

The normal distribution is the famous ``bell curve'':
\[
f(x) = \frac{1}{\sigma\sqrt{2\pi}} e^{-\frac{(x-\mu)^2}{2\sigma^2}}
\]

where $\mu$ is the mean and $\sigma$ is the standard deviation.

We write $X \sim \Normal(\mu, \sigma^2)$ to say ``$X$ follows a normal
distribution with mean $\mu$ and variance $\sigma^2$.''

\begin{warning}
Financial returns are NOT normally distributed! They have ``fat tails''---
extreme events happen much more often than the normal distribution predicts.
We explore this in Chapter~\ref{ch:moments}.
\end{warning}

\subsubsection{Uniform Distribution}

Every outcome in a range is equally likely:
\[
f(x) = \frac{1}{b - a} \quad \text{for } a \leq x \leq b
\]

\subsection{Expected Value and Variance}

\marginnote{%
\textbf{Expected Value}\\[0.5em]
$\E[X] = \displaystyle\sum_x x \cdot P(X = x)$\\[0.5em]
(discrete)\\[1em]
$\E[X] = \displaystyle\int_{-\infty}^{\infty} x \cdot f(x) \, dx$\\[0.5em]
(continuous)
}

The \textbf{expected value} $\E[X]$ is the average outcome you'd expect over
many trials:

For a discrete random variable:
\[
\E[X] = \sum_x x \cdot P(X = x)
\]

Key properties:
\begin{itemize}
    \item $\E[aX + b] = a\E[X] + b$ (linearity)
    \item $\E[X + Y] = \E[X] + \E[Y]$ (additivity)
    \item $\Var(X) = \E[X^2] - (\E[X])^2$
\end{itemize}

\section{Financial Concepts}

Before diving into quantitative methods, let's establish some fundamental
financial concepts.

\subsection{What is a Stock?}

A \textbf{stock} (or share) represents partial ownership of a company. When
you buy a stock, you own a small piece of that company and are entitled to
a share of its profits.

\marginnote{%
\textbf{Stock Price}\\[0.3em]
The price at which buyers and sellers agree to trade ownership.
}

The \textbf{stock price} is determined by supply and demand in the market.
It reflects what buyers are willing to pay and sellers are willing to accept.

\subsection{OHLCV Data}

Market data is typically recorded as \textbf{OHLCV}:

\begin{itemize}
    \item \textbf{O}pen: First price of the trading period
    \item \textbf{H}igh: Highest price during the period
    \item \textbf{L}ow: Lowest price during the period
    \item \textbf{C}lose: Last price of the period
    \item \textbf{V}olume: Number of shares traded
\end{itemize}

This data captures the essential price action within any time frame (day,
hour, minute).

\subsection{Returns}

\marginnote{%
\textbf{Simple Return}\\[0.5em]
$R_t = \dfrac{P_t - P_{t-1}}{P_{t-1}}$\\[1em]
\textbf{Log Return}\\[0.5em]
$r_t = \ln\left(\dfrac{P_t}{P_{t-1}}\right)$
}

A \textbf{return} measures how much an investment gained or lost:

\textbf{Simple return}:
\[
R_t = \frac{P_t - P_{t-1}}{P_{t-1}}
\]

\textbf{Log return} (continuously compounded):
\[
r_t = \ln\left(\frac{P_t}{P_{t-1}}\right)
\]

Log returns have a useful property: they add up over time.
\[
r_{t_1 \to t_n} = r_{t_1} + r_{t_2} + \cdots + r_{t_n}
\]

\subsection{Risk and Volatility}

\textbf{Risk} in finance is typically measured by \textbf{volatility}---the
standard deviation of returns.

\marginnote{%
\textbf{Volatility}\\[0.5em]
$\sigma = \text{std}(r_t)$\\[0.5em]
Often annualized: $\sigma_{\text{annual}} = \sigma_{\text{daily}} \times \sqrt{252}$
}

Higher volatility means more uncertainty---prices swing more wildly.

\textbf{Annualized volatility} scales daily volatility to a yearly figure:
\[
\sigma_{\text{annual}} = \sigma_{\text{daily}} \times \sqrt{252}
\]

(252 is the typical number of trading days per year.)

\subsection{The Risk-Return Tradeoff}

\begin{keyinsight}
Higher expected returns generally require accepting higher risk. There is no
free lunch in finance---if someone promises high returns with low risk, be
skeptical.
\end{keyinsight}

The \textbf{Sharpe ratio} measures risk-adjusted return:
\[
\sharpe = \frac{\E[R] - R_f}{\sigma}
\]

where $R_f$ is the risk-free rate (e.g., Treasury bills).

\section{Programming in Rust}

This book uses Rust for all implementations. Here's a quick introduction for
readers unfamiliar with the language.

\subsection{Why Rust?}

\begin{itemize}
    \item \textbf{Performance}: As fast as C/C++, suitable for
      high-frequency trading
    \item \textbf{Safety}: The compiler catches bugs before runtime
    \item \textbf{Modern}: Great tooling, package management, and community
    \item \textbf{Learning}: Writing correct Rust forces you to think
      carefully---great for learning
\end{itemize}

\subsection{Basic Syntax}

\begin{lstlisting}[language=Rust]
// Variables (immutable by default)
let x = 5;
let mut y = 10;  // mutable
y = 20;

// Functions
fn add(a: i32, b: i32) -> i32 {
    a + b  // no semicolon = return value
}

// Structs (like classes)
struct Stock {
    symbol: String,
    price: f64,
}

// Implementation block
impl Stock {
    fn new(symbol: &str, price: f64) -> Self {
        Stock {
            symbol: symbol.to_string(),
            price,
        }
    }

    fn value(&self, shares: u32) -> f64 {
        self.price * shares as f64
    }
}

// Using the struct
let aapl = Stock::new("AAPL", 150.0);
let portfolio_value = aapl.value(100);
\end{lstlisting}

\subsection{Traits (Interfaces)}

\marginnote{%
\textbf{Traits}\\[0.3em]
Rust's way of defining shared behavior. Similar to interfaces in Java or
protocols in Swift.
}

Traits define shared behavior:

\begin{lstlisting}[language=Rust]
// Define a trait
trait Analyzer {
    fn analyze(&self, data: &[f64]) -> f64;
}

// Implement for a struct
struct MeanAnalyzer;

impl Analyzer for MeanAnalyzer {
    fn analyze(&self, data: &[f64]) -> f64 {
        data.iter().sum::<f64>() / data.len() as f64
    }
}

// Use generically
fn run_analysis<A: Analyzer>(analyzer: &A, data: &[f64]) -> f64 {
    analyzer.analyze(data)
}
\end{lstlisting}

This book uses traits extensively to create modular, reusable components.

\subsection{Error Handling}

\begin{lstlisting}[language=Rust]
// Result type for operations that can fail
fn divide(a: f64, b: f64) -> Result<f64, String> {
    if b == 0.0 {
        Err("Division by zero".to_string())
    } else {
        Ok(a / b)
    }
}

// Using Result
match divide(10.0, 2.0) {
    Ok(result) => println!("Result: {}", result),
    Err(e) => println!("Error: {}", e),
}

// Or use the ? operator
fn calculate() -> Result<f64, String> {
    let x = divide(10.0, 2.0)?;  // propagates error
    Ok(x * 2.0)
}
\end{lstlisting}

\subsection{Iterators}

Rust's iterators are powerful and efficient:

\begin{lstlisting}[language=Rust]
let prices = vec![100.0, 102.0, 101.0, 105.0];

// Map and collect
let returns: Vec<f64> = prices.windows(2)
    .map(|w| (w[1] - w[0]) / w[0])
    .collect();

// Filter
let positive: Vec<f64> = returns.iter()
    .filter(|&&r| r > 0.0)
    .copied()
    .collect();

// Fold (reduce)
let sum: f64 = prices.iter().sum();
let product: f64 = prices.iter().product();
\end{lstlisting}

\section{Running the Examples}

All code examples in this book are available in the companion repository.

\subsection{Prerequisites}

\begin{enumerate}
    \item Install Rust: \url{https://rustup.rs/}
    \item Clone the repository:
\begin{lstlisting}[language=bash]
git clone https://github.com/bitscrafts/quant-lab-rust.git
cd quant-lab-rust
\end{lstlisting}
\end{enumerate}

\subsection{Running Examples}

Each chapter has a corresponding crate with examples:

\begin{lstlisting}[language=bash]
# Chapter 1: Credit Card Fraud
cargo run -p qf-01-fraud --example fraud_analysis

# Chapter 2: Loan Default
cargo run -p qf-02-loan --example loan_analysis

# Chapter 3: Stock Prices
cargo run -p qf-03-stocks --example stock_analysis
\end{lstlisting}

\subsection{Running Tests}

\begin{lstlisting}[language=bash]
# Run all tests
cargo test

# Run tests for a specific crate
cargo test -p qf-01-fraud
\end{lstlisting}

\section{Summary}

This chapter covered:

\begin{itemize}
    \item \textbf{Mathematical notation}: Summation, products, functions,
      logarithms
    \item \textbf{Statistics}: Mean, variance, standard deviation,
      distributions
    \item \textbf{Financial concepts}: Stocks, OHLCV, returns, volatility,
      Sharpe ratio
    \item \textbf{Rust programming}: Variables, functions, structs, traits,
      error handling
\end{itemize}

\vspace{1em}
\noindent\textbf{Next chapter}: We apply these foundations to detect credit
card fraud using anomaly detection.
